\documentclass[10pt,a4paper]{amsart}
\usepackage[margin=19mm]{geometry}
\usepackage{amsmath,amssymb,mathtools}
\usepackage[scr=rsfs,cal=euler]{mathalfa}
\usepackage{xfrac}
\usepackage[unicode,hidelinks,bookmarks=false]{hyperref}
\newcommand{\ie}{i.e.\ }
\newcommand{\cf}{cf.\ }
\newcommand{\resp}{respectively\ }
\newcommand{\Subsection}{\subsection}
\providecommand{\nobreakdash}{\nobreak\hbox{-}}
\setlength{\emergencystretch}{2em}
\multlinegap0pt
\usepackage{bm}
\usepackage{bbm}
\usepackage{dsfont}
\usepackage{xspace}
\usepackage{cancel}
\usepackage{mathtools}
\usepackage[shortlabels]{enumitem}
\usepackage{amsthm}
\makeatletter
\@ifundefined{claim}{\newtheorem{claim}{Claim}}{}
\@ifundefined{claimproof}{\newtheorem{claimproof}{Claimproof}}{}
\@ifundefined{theorem}{\newtheorem{theorem}{Theorem}}{}
\@ifundefined{lemma}{\newtheorem{lemma}[theorem]{Lemma}}{}
\makeatother
\makeatletter
\newcommand{\Aset}{\mathcal A}
\newcommand{\Bset}{\mathcal B}
\newcommand{\Lup}{L^{\mathrm{up}}}
\newcommand{\R}{\mathbb{R}}
\newcommand{\Sset}{\mathcal S}
\newcommand{\T}{\mathsf{T}}
\newcommand{\Tset}{\mathcal T}
\expandafter\ifx\csname claim\endcsname\relax
\newenvironment{claim}{\begin{innerclaim}}{\end{innerclaim}}
\fi
\expandafter\ifx\csname claimproof\endcsname\relax
\newenvironment{claimproof}
  {\par\smallskip\noindent\textit{Proof of the claim.}\ }
\fi
\newcommand{\eps}{\varepsilon}
\newcommand{\ones}{\mathbf{1}}
\DeclareMathOperator{\spn}{span}
\makeatother
\title[Spectral Bounds and Shifted Complexes: Eigenvalues of the Up]{Spectral Bounds and Shifted Complexes: Eigenvalues of the Up-Laplacian via Face Degrees}
\author{Vinayak Gupta}
\date{Source: arXiv:2608.01694v2 (CC BY 4.0)}
\begin{document}
\maketitle

\noindent\emph{Source and licence.} Spectral Bounds and Shifted Complexes: Eigenvalues of the Up-Laplacian via Face Degrees. Version arXiv:2608.01694v2 (CC BY 4.0), distributed under the Creative Commons Attribution 4.0 International licence, as recorded in the arXiv licence metadata for this exact version and shown on the abstract page https://arxiv.org/abs/2608.01694v2. Full rights notice: licenses/arxiv-2608.01694-CC-BY-4.0.txt.\par\smallskip

\noindent\textit{Editorial scope.} Counted unit: the Theorem whose source label is \texttt{thm:main} in the article (source0.tex lines 1772--1778, source label \texttt{thm:main}), delivered together with the article's own complete original proof of it (source0.tex lines 1780--2197, one \texttt{proof} environment of 418 lines). No mathematics is written, repaired or re-derived: statement and proof are the article's own text line for line. Required context retained uncounted: lem:deletion (lines 330--336); lem:gram (lines 360--367); it:gram-b (lines 362--366); lem:compress (lines 371--384); it:comp-a (lines 378--383); lem:switch (lines 413--430); lem:sign (lines 488--494); lem:adj (lines 542--574); it:adj-a (lines 558--573); it:adj-b (lines 558--573); it:adj-c (lines 558--573). Every definition, display and label of those lines is retained unchanged. Omitted from this package and not redistributed: 1--329, 337--359, 367--370, 384--412, 431--487, 495--541, 574--1771, 1779--1779, 2198--2351. Nothing inside a retained excerpt is omitted or altered: every retained body is delivered exactly as the article's lines are written, so no omission receipt of any kind accompanies this package. Citations: the retained excerpts carry no citation command at all, so no bibliography entry of the article is transcribed and no reference is redistributed. Reference adaptation: none was needed and none was made; every reference inside the delivered unit resolves inside it, so no unresolved reference or citation is produced. Printed slips kept literal: no printed word, symbol, hypothesis or number of the counted statement or of its proof was altered. Layout and class: the article's own class is \texttt{amsart} with packages including fontenc, inputenc, amsmath, amssymb, amsthm, mathtools, geometry, enumitem, microtype, xcolor, hyperref; the wrapper is the lane's pinned \texttt{amsart} editorial wrapper at 10pt on A4, which restates the theorem-like environments the retained text needs and adds the article's own macro definitions that the retained text needs (\texttt{\textbackslash Aset}, \texttt{\textbackslash Bset}, \texttt{\textbackslash Lup}, \texttt{\textbackslash R}, \texttt{\textbackslash Sset}, \texttt{\textbackslash T}, \texttt{\textbackslash Tset}, \texttt{\textbackslash eps}, \texttt{\textbackslash ones}, \texttt{\textbackslash spn}), with extra wrapper packages (bm, bbm, dsfont, xspace, cancel, mathtools, enumitem). The delivered numbering is the unit's own. Assets: no retained span contains a figure environment, a graphics-inclusion command, a PDF-page inclusion, a table or a TikZ picture; no image file is copied into this package, the package declares no asset dependency and no image file is redistributed with it. Licence: version arXiv:2608.01694v2 of this article is distributed under the Creative Commons Attribution 4.0 International licence, as recorded in the arXiv licence metadata for this exact version and shown on the abstract page https://arxiv.org/abs/2608.01694v2. A full rights notice is delivered at licenses/arxiv-2608.01694-CC-BY-4.0.txt. Characters: every retained excerpt is pure ASCII, so the delivered engine, XeLaTeX with its default Latin Modern text font, reports no missing character.
\begin{lemma}\label{lem:deletion}
Suppose that $H$ is obtained from $K$ by deleting some $k$-faces while retaining the same row
set for the boundary matrix. Then, for every index $i$,
\[
  \lambda_i\bigl(\Lup_{k-1}(K)\bigr)\ \ge\ \lambda_i\bigl(\Lup_{k-1}(H)\bigr) .
\]
\end{lemma}

\begin{lemma}\label{lem:gram}
Let $B$ be a real matrix.
\begin{enumerate}[label=\textup{(\alph*)},leftmargin=2.6em]
\item The matrices $BB^{\T}$ and $B^{\T}B$ have the same nonzero eigenvalues,
with the same multiplicities.
\item\label{it:gram-b} If $\lambda_2(B^{\T}B)>0$, then $\lambda_2(BB^{\T})=\lambda_2(B^{\T}B)$.
\end{enumerate}
\end{lemma}

\begin{lemma}\label{lem:compress}
Let $A$ be a real symmetric $n\times n$ matrix and let $X\in\R^{n\times r}$ satisfy
$X^{\T}X=I_r$. Then
\[
  \lambda_i(A)\ \ge\ \lambda_i\bigl(X^{\T}AX\bigr),\qquad 1\le i\le r .
\]
In particular:
\begin{enumerate}[label=\textup{(\alph*)},leftmargin=2.6em]
\item\label{it:comp-a} if there is a two-dimensional subspace $U$ with
$z^{\T}Az/z^{\T}z\ge\alpha$ for every nonzero $z\in U$, then $\lambda_2(A)\ge\alpha$;
\item\label{it:comp-b} if $\Aset$ is a set of indices, then
$\lambda_i(A)\ge\lambda_i(A[\Aset])$ for $1\le i\le|\Aset|$.
\end{enumerate}
\end{lemma}

\begin{lemma}\label{lem:switch}
Let $\alpha$ be a ridge and let $\Aset$ be a nonempty collection of facets containing
$\alpha$. The facets can be reoriented so that
\[
  \eps(F,F')=1\qquad\text{for all distinct }F,F'\in\Aset.
\]
Equivalently, there is a diagonal matrix $E$, indexed by the facets and having diagonal entries
in $\{1,-1\}$, such that the down-Laplacian matrix after this reorientation is $EME$ and
\[
  (EME)[\Aset]=kI+J .
\]
Here $E_{F,F}=-1$ means that the orientation of $F$ is reversed, whereas $E_{F,F}=1$ means that
it is left unchanged. Moreover, $EME$ has the same eigenvalues as $M$.

Moreover, let $\beta$ be another ridge and let $\Bset$ be a collection of facets containing
$\beta$. If $|\Aset\cap\Bset|\le1$, then the orientations can be chosen so that both
$(EME)[\Aset]=kI+J$ and $(EME)[\Bset]=kI+J$ hold simultaneously.
\end{lemma}

\begin{lemma}\label{lem:sign}
Let $A,B,C$ be three distinct $k$-faces which are facets of the same abstract
$(k+1)$-simplex. Then, for arbitrary orientations of $A,B,C$,
\[
  \eps(A,B)\,\eps(B,C)\,\eps(C,A)=-1 .
\]
\end{lemma}

\begin{lemma}\label{lem:adj}
Let $\sigma\ne\tau$ be $(k-1)$-faces, put
\[
  \rho=\sigma\cap\tau,\qquad q=|\sigma\setminus\tau|=|\tau\setminus\sigma|\ge1,
  \qquad |\rho|=k-q ,
\]
and let $F=\sigma\cup\{u\}$ and $G=\tau\cup\{v\}$ be $k$-faces with $F\ne G$, where
$u\notin\sigma$ and $v\notin\tau$. Then
\[
  |F\cap G|=
  \begin{cases}
    k-q+1, & u=v,\\[2pt]
    k-q+|\sigma\cap\{v\}|+|\tau\cap\{u\}|, & u\ne v.
  \end{cases}
\]
Consequently $|F\cap G|\le k-q+2$, and:
\begin{enumerate}[label=\textup{(\alph*)},leftmargin=2.6em]
\item\label{it:adj-a} if $q\ge3$, then $F$ and $G$ never share a ridge;
\item\label{it:adj-b} if $q=1$, write $\sigma=\rho\cup\{a\}$ and $\tau=\rho\cup\{b\}$. If in
addition $u\ne b$ and $v\ne a$, then $F$ and $G$ share a ridge if and only if $u=v$, and in
that case $F\cap G=\rho\cup\{u\}$ and $u\notin\rho\cup\{a,b\}$;
\item\label{it:adj-c} if $q=2$, write $\sigma=\rho\cup\{a_1,a_2\}$ and
$\tau=\rho\cup\{b_1,b_2\}$. Then $F$ and $G$ share a ridge if and only if
\[
  F=F_i:=\sigma\cup\{b_i\}\quad\text{and}\quad G=G_j:=\tau\cup\{a_j\}
\]
for some $i,j\in\{1,2\}$. Moreover,
$W=\sigma\cup\tau=\rho\cup\{a_1,a_2,b_1,b_2\}$ is a $(k+2)$-element vertex set, not
necessarily a face of $K$, and $F_1,F_2,G_1,G_2$ are four pairwise distinct
$(k+1)$-element subsets of $W$. Furthermore,
$F_i\cap G_j=\rho\cup\{a_j,b_i\}$ is a ridge for \emph{every} pair $(i,j)$.
\end{enumerate}
\end{lemma}

\begin{theorem}\label{thm:main}
Let $K$ be a finite $k$-dimensional simplicial complex, where $k\ge1$, with at least two
$k$-faces, and let $\Lup_{k-1}(K)$ be its unnormalized $(k-1)$-dimensional up-Laplacian. Then
\[
  \lambda_2\bigl(\Lup_{k-1}(K)\bigr)\ \ge\ d_2(K)+k-1 .
\]
\end{theorem}

\begin{proof}
Choose distinct $(k-1)$-faces $\sigma$ and $\tau$ such that
\[
  d(\sigma)=d_1(K),\qquad d(\tau)=d_2(K),
\]
and put $d:=d_2(K)$. Such a choice is possible because $K_{k-1}$ has at least two elements, so
that $d_1(K)$ and $d_2(K)$ are the two largest members of a list of at least two upper degrees.
We first assume that $d\ge2$. The case $d=1$ will be treated at the end. 

\medskip
\noindent

Retain \emph{all} $d$ facets containing $\tau$, and retain exactly $d$ facets containing
$\sigma$; the latter is possible because $d(\sigma)=d_1(K)\ge d_2(K)=d$. If $\sigma$ and $\tau$
lie in a common $k$-face, include that common face among the retained facets containing
$\sigma$. Let $H$ be the subcomplex consisting of the entire $(k-1)$-skeleton of $K$, together
with the retained $k$-faces. Define
\[
  \Sset=\{F\in H_k:\ \sigma\subset F\},\qquad \Tset=\{F\in H_k:\ \tau\subset F\} .
\]
Thus $\Sset$ and $\Tset$ are the two selected facet stars, and
\[
  |\Sset|=|\Tset|=d .
\]
Because two distinct $(k-1)$-faces can lie together in at most one $k$-face, a facet lying in
both stars would contain $\sigma\cup\tau$ and hence be uniquely determined; therefore
\[
  |\Sset\cap\Tset|\le1 .
\]
After the reduction, the two possible configurations are exactly
\[
  \Sset\cap\Tset=\varnothing \qquad\text{or}\qquad \Sset\cap\Tset=\{F_0\},
\]
where $F_0$ is the common $k$-face in $\Sset$ and $\Tset$. Since $H$ has the same $(k-1)$-faces
as $K$, Lemma~\ref{lem:deletion} applies and gives
\[
  \lambda_2\bigl(\Lup_{k-1}(K)\bigr)\ \ge\ \lambda_2\bigl(\Lup_{k-1}(H)\bigr) .
\]
It is therefore enough to establish the required estimate for $H$.

\medskip
\noindent

Consider the down Gram matrix
\[
  M={B_k(H)}^{\T}B_k(H) .
\]
 $M$ is
indexed by the $2d$ or $2d-1$ retained facets, and its diagonal is constant. Indeed every
diagonal entry of $M$ is $k+1$, so we may write
\[
  M=kI+N ,
\]
where $N$ has unit diagonal and $N_{F,G}=\eps(F,G)$ for adjacent $F\ne G$, and $N_{F,G}=0$
otherwise. Every eigenvalue of $M$ exceeds the corresponding eigenvalue of $N$ by exactly $k$,
so it suffices to bound $\lambda_2(N)$ from below by $d-1$.

Since $|\Sset\cap\Tset|\le1$, Lemma~\ref{lem:switch} applies to the pair of stars
$\Sset,\Tset$, associated with the $(k-1)$-faces $\sigma,\tau$. We may therefore replace $M$ by
$EME$ and assume from now on that
\[
  N[\Sset]=J_d,\qquad N[\Tset]=J_d .
\]
The switching does not change the spectrum, and by Lemma~\ref{lem:switch} the switched matrix
is again a down-Laplacian matrix of $H$; consequently the incidence
products $\eps$ computed in the new orientations still satisfy Lemma~\ref{lem:sign}, which we
shall use repeatedly. We now treat the two possible intersections of the selected stars
separately.

\medskip
\noindent\textbf{Case 1. $\Sset\cap\Tset=\varnothing$.}
Order the retained facets with those in $\Sset$ first and those in $\Tset$ second. Then
\[
  N=
  \begin{pmatrix}
    J_d & C\\
    C^{\T} & J_d
  \end{pmatrix},
  \qquad
  C=N[\Sset,\Tset] .
\]
Define
\[
  x=\begin{pmatrix}\ones_d\\\mathbf{0}_d\end{pmatrix},
  \qquad
  y=\begin{pmatrix}\mathbf{0}_d\\\ones_d\end{pmatrix},
  \qquad
  s=\ones^{\T}C\ones
   =\sum_{F\in\Sset}\sum_{G\in\Tset}N_{F,G} .
\]
Thus $x$ and $y$ are the $0$--$1$ vectors indicating membership in $\Sset$ and $\Tset$,
respectively. Since the stars are disjoint, every pair $(F,G)\in\Sset\times\Tset$ has
$F\ne G$, so that $C_{F,G}=\eps(F,G)$ when $F$ and $G$ are adjacent and $C_{F,G}=0$ otherwise.

We first exclude one exceptional case.

\begin{claim}\label{cl:pivot}
If $\Sset\cap\Tset=\varnothing$, then $\sigma\cup\tau$ is not a $k$-face of $K$.
\end{claim}

\begin{claimproof}
If $|\sigma\cup\tau|\ne k+1$ there is nothing to prove, since a $k$-face has exactly $k+1$
vertices. So suppose $|\sigma\cup\tau|=k+1$ and that $F:=\sigma\cup\tau$ is a $k$-face of
$K$. Then $\tau\subset F$, and the definition of $H$ retains every $k$-face containing $\tau$;
hence $F\in\Tset$. Also $\sigma\subset F$, so $F$ is a common $k$-face of $\sigma$ and $\tau$.
By the definition of $H$, this common face is also retained among the facets containing
$\sigma$; hence $F\in\Sset$. This contradicts
$\Sset\cap\Tset=\varnothing$.
\end{claimproof}

\begin{claim}\label{cl:s}
In the disjoint-star case, $|s|\le d$.
\end{claim}

\begin{claimproof}
Put
\[
  q=|\sigma\setminus\tau|=|\tau\setminus\sigma| .
\]
Since $\sigma\ne\tau$ and the two faces have the same dimension, $q\ge1$. We divide the
argument according to the value of $q$; the three subcases $q\ge3$, $q=1$ and $q=2$ are
exhaustive and mutually exclusive.

\smallskip
\noindent\emph{Subcase 1.1: $q\ge3$.} Every facet in $\Sset$ is obtained from $\sigma$ by
adding one vertex, and every facet in $\Tset$ is obtained from $\tau$ by adding one vertex.
Since $\sigma$ and $\tau$ differ in at least three vertices, Lemma~\ref{lem:adj}\ref{it:adj-a}
shows that no such pair can share a ridge. Hence $C$ is the zero matrix, and therefore $s=0$.

\smallskip
\noindent\emph{Subcase 1.2: $q=1$.} Write
\[
  \sigma=\rho\cup\{a\},\qquad \tau=\rho\cup\{b\},\qquad |\rho|=k-1 .
\]
Every facet in $\Sset$ has the form $F_u=\rho\cup\{a,u\}$ with $u\notin\sigma$, and every facet
in $\Tset$ has the form $G_v=\rho\cup\{b,v\}$ with $v\notin\tau$. We first check that the
degenerate values $u=b$ and $v=a$ do not occur. If $u=b$ then $F_u=\rho\cup\{a,b\}=\sigma\cup\tau$
would be a $k$-face of $K$, contradicting Claim~\ref{cl:pivot}; the same argument excludes
$v=a$.

With $u\ne b$ and $v\ne a$ available, Lemma~\ref{lem:adj}\ref{it:adj-b} shows that $F_u$ and
$G_v$ share a ridge precisely when $u=v$. Thus every row and every column of $C$ contains at
most one nonzero entry. Consequently $C$ has at most $d$ nonzero entries, each equal to $1$ or
$-1$, and hence
\[
  |s|\le d .
\]

\smallskip
\noindent\emph{Subcase 1.3: $q=2$.} Write
\[
  \sigma=\rho\cup\{a_1,a_2\},\qquad \tau=\rho\cup\{b_1,b_2\},
\]
where the four displayed vertices are distinct. By Lemma~\ref{lem:adj}\ref{it:adj-c} the only
facets that can participate in a cross-interaction are
\[
  F_i=\sigma\cup\{b_i\},\qquad G_j=\tau\cup\{a_j\},\qquad i,j\in\{1,2\},
\]
These four are $(k+1)$-element subsets of the same $(k+2)$-element vertex set
$W=\sigma\cup\tau$, which need not be a face of $K$, and for every $i,j$,
\[
  F_i\cap G_j=\rho\cup\{b_i,a_j\},
\]
which is a ridge. Thus every retained $F_i$ interacts with every retained $G_j$.

Assume that both $G_1$ and $G_2$ are retained. For each retained $F_i$, the three facets
$F_i,G_1,G_2$ are distinct $(k+1)$-element subsets of the same vertex set $W$. By the chosen
diagonal switching, the $\Tset$-block of $N$ is positive. Since $G_1,G_2\in\Tset$, we have
\[
  \eps(G_1,G_2)=N_{G_1,G_2}=1 .
\]
Lemma~\ref{lem:sign} applied to the triple $F_i,G_1,G_2$ gives
$\eps(F_i,G_1)\eps(G_1,G_2)\eps(G_2,F_i)=-1$, whence
\[
  \eps(F_i,G_1)=-\eps(F_i,G_2) .
\]
Thus, for every retained $F_i$, the two entries corresponding to $G_1$ and $G_2$ have opposite
signs and hence sum to zero. Similarly, if both $F_1$ and $F_2$ are retained, the two entries
in each retained column sum to zero, because $F_1,F_2\in\Sset$ gives
$\eps(F_1,F_2)=1$ and Lemma~\ref{lem:sign} applied to $F_1,F_2,G_j$ gives
$\eps(F_1,G_j)=-\eps(F_2,G_j)$.

It remains to combine these observations. If neither $F_1$ nor $F_2$ is retained, or neither
$G_1$ nor $G_2$ is retained, then the cross block is zero and $s=0$. If both $G_1,G_2$ are
retained, its rows cancel; if both $F_1,F_2$ are retained, its columns cancel. Otherwise, at
most one $F_i$ and at most one $G_j$ are retained, so there is at most one nonzero cross-entry
and $|s|\le1$. In every configuration,
\[
  |s|\le1\le d .
\]
This completes all three subcases and proves the claim.
\end{claimproof}

We now use Claim~\ref{cl:s} to obtain the required eigenvalue estimate. From the definitions
of $x$ and $y$,
\[
  x^{\T}x=y^{\T}y=d,\qquad x^{\T}y=0 .
\]
The block form of $N$ gives
\[
  x^{\T}Nx=y^{\T}Ny=d^2,
  \qquad
  x^{\T}Ny=s .
\]
Set
\[
  Q=\begin{pmatrix}x/\sqrt d & y/\sqrt d\end{pmatrix} .
\]
The columns of $Q$ are orthonormal, so $Q^{\T}Q=I_2$. The preceding identities give
\[
  Q^{\T}NQ
  =\begin{pmatrix}
      d & s/d\\[2pt]
      s/d & d
    \end{pmatrix} .
\]
The smaller eigenvalue of this $2\times2$ matrix is
\[
  d-\frac{|s|}{d}\ \ge\ d-1
\]
by Claim~\ref{cl:s}. Applying Lemma~\ref{lem:compress} gives
\[
  \lambda_2(N)\ \ge\ d-1 ,
\]
and hence, since $M=kI+N$,
\[
  \lambda_2(M)\ \ge\ k+d-1 .
\]

\medskip
\noindent\textbf{Case 2. $\Sset\cap\Tset=\{F_0\}$.}
Because $F_0$ contains both $\sigma$ and $\tau$, we have $|\sigma\cup\tau|\le k+1$, so these two
$(k-1)$-faces differ in exactly one vertex, that is $q=1$. Write
\[
  \rho=\sigma\cap\tau,\qquad \sigma=\rho\cup\{a\},\qquad \tau=\rho\cup\{b\} .
\]
Then $\sigma\cup\tau$ has $k+1$ vertices and is contained in $F_0$, so
\[
  F_0=\rho\cup\{a,b\} .
\]
Every other facet in $\Sset$ has the form $F_u=\rho\cup\{a,u\}$ with $u\notin\sigma$ and, since
$F_u\ne F_0$, with $u\ne b$; every other facet in $\Tset$ has the form $G_v=\rho\cup\{b,v\}$
with $v\notin\tau$ and $v\ne a$. The next claim describes all entries of $N$ between
$\Sset\setminus\{F_0\}$ and $\Tset\setminus\{F_0\}$ and gives the resulting block form of $N$.

\begin{claim}\label{cl:structure}
The cross-interactions between $\Sset\setminus\{F_0\}$ and $\Tset\setminus\{F_0\}$ form a
partial matching, and every matched cross-entry is $-1$. Consequently, with the ordering
\[
  F_0,\quad \Sset\setminus\{F_0\},\quad \Tset\setminus\{F_0\},
\]
the matrix $N$ has the form
\[
  N=\begin{pmatrix}
    1 & \ones^{\T} & \ones^{\T}\\
    \ones & J_{d-1} & -P\\
    \ones & -P^{\T} & J_{d-1}
  \end{pmatrix},
\]
where $P$ is a $0,1$ partial permutation matrix.
\end{claim}

\begin{claimproof}
Since $u\ne b$ and $v\ne a$, Lemma~\ref{lem:adj}\ref{it:adj-b} applies and shows that $F_u$ and
$G_v$ share a ridge if and only if $u=v$. Concretely, when $u=v$ the intersection
$F_u\cap G_v=\rho\cup\{u\}$ has $k$ vertices, whereas for $u\ne v$ their intersection is $\rho$,
which has only $k-1$ vertices. Hence the cross-interactions form a partial matching, and the
corresponding block has at most one nonzero entry in each row and each column, every such entry
being either $1$ or $-1$. It remains to determine their signs.

Consider any matched pair $F_u,G_u$. Since $u\notin\rho\cup\{a,b\}$, the set
\[
  W_u=\rho\cup\{a,b,u\}
\]
has $k+2$ vertices, and
\[
  F_0=W_u\setminus\{u\},\qquad
  F_u=W_u\setminus\{b\},\qquad
  G_u=W_u\setminus\{a\} .
\]
Thus $F_0,F_u,G_u$ are three distinct $(k+1)$-element subsets of $W_u$, so
Lemma~\ref{lem:sign} applies to them. Moreover, the switched star blocks are positive. Since
$F_0,F_u\in\Sset$ and $F_0,G_u\in\Tset$, we have
\[
  \eps(F_0,F_u)=N_{F_0,F_u}=1,\qquad \eps(G_u,F_0)=N_{G_u,F_0}=1 .
\]
Lemma~\ref{lem:sign} therefore gives
\[
  \eps(F_0,F_u)\,\eps(F_u,G_u)\,\eps(G_u,F_0)=-1 .
\]
Substituting the two known signs yields $\eps(F_u,G_u)=-1$. Since the matched pair was
arbitrary, every nonzero cross-entry is $-1$. Hence the cross block is $-P$, where $P$ is a
$0$--$1$ partial permutation matrix of order $d-1$.

For the displayed shape of $N$ it remains to identify the entries in the row and column of
$F_0$. If $F\in\Sset\setminus\{F_0\}$ then $F$ and $F_0$ both lie in $\Sset$, so
$N_{F_0,F}=1$; if $G\in\Tset\setminus\{F_0\}$ then $G$ and $F_0$ both lie in $\Tset$, so
$N_{F_0,G}=1$. The diagonal entry is $N_{F_0,F_0}=1$, and the two diagonal blocks are $J_{d-1}$
because $\Sset$ and $\Tset$ have positive blocks. This proves both assertions of the claim.
\end{claimproof}

With the ordering in Claim~\ref{cl:structure}, define
\[
  x=\begin{pmatrix}1\\\ones_{d-1}\\\mathbf{0}_{d-1}\end{pmatrix},
  \qquad
  y=\begin{pmatrix}1\\\mathbf{0}_{d-1}\\\ones_{d-1}\end{pmatrix} .
\]
Thus the coordinate of $x$ indexed by a facet is $1$ exactly when that facet lies in $\Sset$;
the analogous statement holds for $y$ and $\Tset$.

Let $r$ be the number of nonzero entries of $P$. Since $P$ is a partial permutation matrix of
order $d-1$,
\[
  0\le r\le d-1 .
\]
 Note that $x$ and $y$ are
linearly independent: they are $0$--$1$ vectors with $d\ge2$ ones each, and they agree in only
one coordinate.

\begin{claim}\label{cl:rayleigh}
Every nonzero vector in $\spn\{x,y\}$ has Rayleigh quotient with respect to $N$ at least $d-1$.
\end{claim}

\begin{claimproof}
Since the two stars meet only in $F_0$,
\[
  x^{\T}x=y^{\T}y=d,\qquad x^{\T}y=1 .
\]
Since $N[\Sset]=N[\Tset]=J_d$,
\[
  x^{\T}Nx=y^{\T}Ny=d^2 .
\]
Since $P$ is a $0$--$1$ matrix with $r$ nonzero entries,
$\ones_{d-1}^{\T}P\ones_{d-1}=r$. Direct multiplication gives
\[
  x^{\T}Ny=1+(d-1)+(d-1)-r=2d-1-r .
\]
Put
\[
  p=x+y,\qquad w=x-y .
\]
These vectors are orthogonal, because
\[
  p^{\T}w=x^{\T}x-y^{\T}y=0 ,
\]
and they are nonzero, since $p^{\T}p=2(d+1)>0$ and $w^{\T}w=2(d-1)>0$, the latter using
$d\ge2$. They are also orthogonal with respect to the quadratic form of $N$, since
\[
  p^{\T}Nw=x^{\T}Nx-x^{\T}Ny+y^{\T}Nx-y^{\T}Ny=d^2-(2d-1-r)+(2d-1-r)-d^2=0 .
\]
Since $\{p,w\}$ is a basis of $\spn\{x,y\}$ and both mixed terms vanish, the Rayleigh quotient
of any nonzero vector $z=\alpha p+\beta w$ is a weighted average of the Rayleigh quotients of
$p$ and $w$. It therefore suffices to show that both of these quotients are at least $d-1$.

For $w=x-y$,
\[
  w^{\T}w=2(d-1),\qquad
  w^{\T}Nw=2\bigl(d^2-(2d-1-r)\bigr)=2\bigl({(d-1)}^{2}+r\bigr) .
\]
Therefore
\[
  \frac{w^{\T}Nw}{w^{\T}w}=\frac{{(d-1)}^{2}+r}{d-1}=d-1+\frac{r}{d-1}\ \ge\ d-1 ,
\]
since $r\ge0$ and $d-1>0$.

For $p=x+y$,
\[
  p^{\T}p=2(d+1),\qquad p^{\T}Np=2\bigl(d^2+2d-1-r\bigr) .
\]
Consequently
\[
  \frac{p^{\T}Np}{p^{\T}p}
  =\frac{d^2+2d-1-r}{d+1}
  =d-1+\frac{2d-r}{d+1} >\ d-1 ,
\]
because $r\le d-1<2d$. Both orthogonal directions therefore have Rayleigh quotient at least
$d-1$, and the mixed quadratic term vanishes. The claim follows.
\end{claimproof}

By Claim~\ref{cl:rayleigh} and Lemma~\ref{lem:compress}\ref{it:comp-a}, applied to the
two-dimensional subspace $\spn\{x,y\}$,
\[
  \lambda_2(N)\ \ge\ d-1 .
\]
Hence, also in the one-common-facet case,
\[
  \lambda_2(M)\ \ge\ k+d-1 .
\]

\medskip
\noindent

For $d\ge2$, the preceding two cases give $\lambda_2(M)\ge k+d-1>0$.
Lemma~\ref{lem:gram}\ref{it:gram-b} and deletion monotonicity therefore give
\[
  \lambda_2\bigl(\Lup_{k-1}(K)\bigr)
  \ \ge\ \lambda_2\bigl(\Lup_{k-1}(H)\bigr)
  \ =\ \lambda_2(M)
  \ \ge\ k+d-1
  \ =\ d_2(K)+k-1 .
\]

\medskip
It remains to prove the inequality when $d=1$. Since $K$ has at least two $k$-faces, choose
distinct $F_1,F_2\in K_k$. The
principal submatrix of $M(K)$ indexed by these faces has the form
\[
  \begin{pmatrix}k+1&\eps\\ \eps&k+1\end{pmatrix},
  \qquad \eps\in\{0,1,-1\}.
\]
Its smaller eigenvalue is $k+1-|\eps|\ge k>0$. Hence interlacing and
Lemma~\ref{lem:gram}\ref{it:gram-b} give
\[
  \lambda_2\bigl(\Lup_{k-1}(K)\bigr)
  =\lambda_2\bigl(M(K)\bigr)\ge k=d+k-1 .
\]
The proof is complete.
\end{proof}

\end{document}
